\documentclass[11pt]{article}
\usepackage{fontspec}
\setmainfont{Latin Modern Roman}
\usepackage{amsmath,amssymb,mathtools,unicode-math}
\setmathfont{Latin Modern Math}
\usepackage{xurl,hyperref}
\hypersetup{hidelinks,pdftitle={Problem statement}}
\newcommand{\sref}[2]{\hyperlink{source:#1:#2}{[S#2]}}
\newcommand{\cataloguescope}[1]{\paragraph{Mathematical question.}#1}
\begin{document}
\section{Research Problems in Function Theory — Problem 7.57}\label{statement}

\subsection{Definitions and mathematical statement}

\paragraph{Shared notation.}$\mathbb N=\{1,2,\ldots\}$, and $\mathbb Z,\mathbb Q,\mathbb R,\mathbb C$ denote the integers, rationals, reals and complex numbers. Any other conventions are specified locally.


Fix $n\ge2$. Put $H_\pm=\{x\in\mathbb R^n:\pm x_n>0\}$ and $\sigma(x_1,\ldots,x_n)=(x_1,\ldots,x_{n-1},-x_n)$. For nonempty compact $E\subset H_+$ and $F\subset H_-$ set $F^*=\sigma(F)$. For sets $A,B$, let $\Delta(A,B)$ be all curves in the one-point compactification $\mathbb R^n\cup\{\infty\}$ joining them. The $n$-modulus is
\[M_n(\Delta)=\inf_\rho\int_{\mathbb R^n}\rho(x)^n\,dx,
\]
where $\rho:\mathbb R^n\to[0,\infty]$ is Borel and $\int_\gamma\rho\,ds\ge1$ for every locally rectifiable joining curve; nonrectifiable curves impose no condition. The infimum is $+\infty$ when there is no admissible density.
\paragraph{Statement recorded in the supplied source.}
Does
\[M_n(\Delta(E,F))\le M_n(\Delta(E,F^*))
\]
hold for every such pair of compact sets? \sref{statement}{2}
\subsection{Short English statement}
Does reflecting the lower of two compact sets across the separating hyperplane increase the modulus of the family of curves joining the sets?
\subsection{Sources}
\begin{description}
\item[\hypertarget{source:statement:1}{S1}] ULAM.AI and the UnsolvedMath Contributors, UnsolvedMath: A Curated Collection of Open Mathematics Problems, catalogue entry. CC BY 4.0. \url{https://huggingface.co/datasets/ulamai/UnsolvedMath}.
\item[\hypertarget{source:statement:2}{S2}] Walter K. Hayman and Edward F. Lingham, Research Problems in Function Theory (2018), Problem 7.57, its complete update and relevant preceding definitions. \url{https://arxiv.org/abs/1809.07200}.
\end{description}
\label{end:statement}
\end{document}
